% !TEX root =  ../Report.tex

\chapter{Introduction}
\section{Introduction}
\label{sec:Introduction}


Combinatorics is a field of study in mathematics that primarily focuses on
the notion of counting objects with certain properties. Over time, this notion
has shifted, especially in the subfield of algebraic combinatorics, where
there is no clear notion of the object that we are counting,
and the numbers express something more abstract \cite{pak_WhatCombinatorialInterpretation_2022}.
Researchers have questioned whether it is possible to address the inverse question,
or more informally, to associate a sequence of numbers to sets of objects
\cite{ikenmeyer_WhatWhatNot_2022,pak_WhatCombinatorialInterpretation_2022,ikenmeyer_PositivitySymmetricGroup_2024}.


To understand this notion, we provide an example from the papers by Ikenmeyer and Pak \cite{pak_WhatCombinatorialInterpretation_2022,ikenmeyer_WhatWhatNot_2022}:
Assume that $f(G)$ counts the number of 3-colourings of a graph $G$. A natural question follows: what does 
$h(G) \doteq f(G)/ 6$ count?
As shown in \cite{pak_WhatCombinatorialInterpretation_2022,ikenmeyer_WhatWhatNot_2022},
for every colouring $\chi$, there are $3!$ additional colourings since
we can permute the colours we assigned. Therefore, we only need to count
one of the permutations, or the smallest colouring. We can therefore argue that $h(G)$ 
counts the smallest possible 3-colourings of a graph $G$.
We denote this as a ``combinatorial interpretation of $h(G)$''.

However not all non-negative functions have such a clear interpretation.
We can utilise an example by Ikenmeyer et al. \cite{ikenmeyer_WhatWhatNot_2022},
where we recall a combinatorial variant of the \textit{Pigeonhole principle},
given a circuit $C:\{0,1\}^n \to \{0,1\}^n$, find $x$ such that $C(x) = 0^n$
or distinct $x,y \in \{0,1\}^n$ such that $C(x)=C(y)$ and $C(C(x)) \neq 0$.
$C$ essentially maps
$\{0,1\}^n \to \{0,1\}^n \smallsetminus \{0^n\}$.
Every point that is mapped to $0^n$ is considered a violation.
The problem counts every pair of points that is mapped to the same bin
or to a point that is a violation. We add the restriction
of $C(C(x)) \neq 0$ to avoid double counting. We can clearly observe that, if $e(C)$
counts the number of solutions, under the pigeonhole principle, it must be the case $e(C) \geq 1$ since we either have a violation
or there exists at least one pair of points that is mapped to the same bin by the pigeonhole principle.
This raises the question: does the function $e'(C) \doteq e(C) -1$
have a combinatorial interpretation?
Ikenmeyer et al. \cite{ikenmeyer_WhatWhatNot_2022} demonstrated that
this problem may not have a combinatorial interpretation.
This leads to the question of how can we formally represent the meaning
of a combinatorial interpretation, and what can we do to verify whether
a function has one?

In recent years, there has been a revolutionary initiative to answer the aforementioned questions,
with the help of counting complexity theory and the complexity class of \textbf{\#P}.
In fact, being able to find such definitions or interpretations is crucial,
as it allows us to utilise tools from enumerative combinatorics to understand
and reveal hidden structures and properties of such functions \cite{pak_WhatCombinatorialInterpretation_2022,ikenmeyer_WhatWhatNot_2022}.
Moreover, there are several problems or numbers, such as \textit{Kronecker coefficients} \cite{makar_AnalysisKroneckerProduct_1949},
whose combinatorial interpretation could lead to groundbreaking breakthroughs
or resolutions of key conjectures such as $\textbf{P} \neq \textbf{NP}$ \cite{ikenmeyer_WhatWhatNot_2022, ikenmeyer_VanishingKroneckerCoefficients_2017}.


In our current work, we focus on extending the work done by
Ikenmeyer et al. \cite{ikenmeyer_WhatWhatNot_2022}, where they focused on the creation of frameworks
that determine whether $f \in^? \textbf{\#P}$, by looking
at the subclasses of $\textbf{\#TFNP}-1$ problems, a class of problems which guarantee a solution.
More specifically, we look into the \textbf{PPAD} class which
seems to have several \textit{complete} problems where under decrementation,
may or may not have a combinatorial interpretation \cite{ikenmeyer_WhatWhatNot_2022}.
This implies that class \textbf{PPAD} acts as some sort of barrier
to $\textbf{\#P}$ which raises questions as to its combinatorial nature.
To accomplish this task,
we studied \textit{PureCircuit} problem
which was created by Deligkas et al. \cite{deligkas_PureCircuitTightInapproximability_2024},
as we hope its \textbf{PPAD}-completeness and circuit-like nature may give us
useful insights to the counting complexity of $\textbf{\#PPAD}-1$.


We will now give an extensive introduction to the literature and related problems
and consequently offer formal and precise objectives that capture our goal.
We will then go over the main results of our theory research with the conjectures and possible
directions that we uncovered. Moreover, we will give an overview of the visualisation software
we developed to verify our theoretical results and give an overview of its capabilities and limitations.
Lastly we will give brief analysis of our project management, our personal reflections
and the future directions.


% In our work we focus specifically on the complexity class of $\textsc{PPAD}$, under the lens of $\text{PureCircuit}$,
% due to its circuit-based nature.
% We aspire to uncover many insights of the $\textsc{\#PureCircuit} -1$ problem
% and explore the boundaries of $\textsc{\#PPAD} -1$.